% sections/memory.tex — Memory Model
\section{Memory Model}
\label{sec:memory}

The lifecycle and K2K calculus of \Cref{sec:model,sec:k2k} are
parameterized by a memory model: the bytes of \(\mathit{snap}\) and
\(\mathit{live}\), the queues \(Q(d)\), the dead-letter store \(X\),
and the idempotency cache \(I\) all live somewhere in the device's
address space.  Where they live—and how the runtime constrains
their visibility, persistence, and access ordering—affects which
small-step traces the actor system can realize.  This section is
the \emph{assumption boundary} between the abstract calculus and
Hopper's actual execution and memory semantics: it establishes
which device behaviors we take as given, which placement decisions
are the runtime's design choices, and which orderings our proofs
rely on.  The three layers are
\textbf{(i)}~hardware-documented facts about Hopper's memory
hierarchy, DSMEM, cooperative groups, and atomics
(\Cref{sec:memory:hierarchy},
\citet{nvidia2022hopper,nvidia2023cooperative,nvidia2023hopperguide});
\textbf{(ii)}~runtime design choices about where queues,
snapshots, and caches live
(\Cref{sec:memory:placement});
and \textbf{(iii)}~proof assumptions about release-acquire
ordering and snapshot atomicity that the K2K and supervisor
theorems depend on
(\Cref{sec:memory:consistency,sec:memory:snapshot}).  We are
explicit about which claim rests on which layer.

\subsection{Hierarchy}
\label{sec:memory:hierarchy}

NVIDIA Hopper exposes a five-tier memory hierarchy:

\begin{enumerate}
\item \textbf{Registers}: per-thread, lifetime equal to the
issuing warp's stay on the SM, no inter-thread visibility without an
explicit shuffle.
\item \textbf{Shared memory (SMEM)}: per-block, on-chip, $\sim$30
cycles of latency, partitioned into 32 banks; visible to all threads
in the block; reset on block exit.
\item \textbf{Distributed shared memory (DSMEM)}: per-cluster,
on-chip but with cross-block addressing via
\texttt{cluster.map\_shared\_rank()}; $\sim$30 cycles of latency
within a GPC~\citep{nvidia2022hopper}; reset on cluster exit.
\item \textbf{L1 / L2 caches}: per-SM L1, per-device L2 ($\sim$50~MiB
on H100); software-transparent, but \texttt{\_\_ldg()} and
\texttt{cuda::memcpy\_async} give the kernel some control over
caching policy.
\item \textbf{HBM (global memory)}: per-device, off-chip, $\sim$400
cycles of latency, persistent across kernel launches.
\end{enumerate}

Persistent actors necessarily live in HBM: they survive across
kernel invocations and the on-chip tiers do not.  But the tiers
above HBM are essential to performance: a hot inbox lives partly in
SMEM (the in-kernel ring), partly in DSMEM (the cross-block
forwarding ring), partly in L2 (recently-touched envelopes), and
only its overflow tail in HBM.  The semantics must allow this
distribution while preserving the abstract calculus.

\subsection{Placement of Model Components}
\label{sec:memory:placement}

We map each component of the formal model to its physical residence
in the runtime:

\begin{itemize}
\item \(\mathit{snap}\) (per-actor snapshot): HBM, in a
supervisor-private region not addressable by the actor.

\item \(\mathit{live}\) (per-actor live state): HBM under the actor's
own ownership; the actor mutates it directly via \textsc{Process}.
The kernel may cache portions in registers and SMEM during a
\textsc{Process} step; the live HBM copy is the canonical version.

\item \(Q(d)\) (inbox queue): tiered. The first \(N_{\mathrm{smem}}\)
slots live in SMEM (only when \(d\) is in the same block as the
sender, i.e., \textsc{Smem} tier); the next \(N_{\mathrm{dsmem}}\)
slots live in DSMEM (when \(d\) is in the same cluster); the
overflow tail lives in HBM.  The runtime chooses
\(N_{\mathrm{smem}}, N_{\mathrm{dsmem}}\) at sub-broker creation
time.

\item \(D\) (delivered set): not materialized at runtime; the model
treats it as a derived quantity. The runtime maintains only the
\emph{count} and per-(source,\,destination) HLC of the most recent
delivery, in HBM.

\item \(X\) (dead-letter store): HBM, a supervisor-managed ring
buffer per sub-broker. Its capacity is configurable; eviction policy
is FIFO.

\item \(I\) (idempotency cache): HBM, a Cuckoo hash with
\(2^{16}\)-entry capacity per sub-broker, evicting on insert collisions.

\item HLC values: registers and SMEM during a \textsc{Process} step;
the canonical per-actor HLC lives in the actor's HBM state vector
(\Cref{def:actor-state}, \(\mathit{hlc}\) field).
\end{itemize}

The dead-letter store and idempotency cache are intentionally HBM-only:
they must survive both \textsc{Process} steps and \textsc{Restart}.
The inbox tier choice is a performance-versus-capacity trade-off
rather than a correctness one; the calculus remains the same
regardless of which tier the slots are physically in.

\subsection{Memory Consistency}
\label{sec:memory:consistency}

The GPU's memory consistency model
(\citet{lustig2014abstraction,hower2014}) is weaker than sequential
consistency: a thread's writes are observable by other threads only
in a partial order constrained by acquire/release fences.  The K2K
calculus relies on three concrete constraints from that model:

\begin{enumerate}
\item \textbf{Producer release}: at the end of \textsc{SendOK},
the producer issues a
\texttt{cuda::\allowbreak atomic\_thread\_fence(\allowbreak
release,\ scope::cluster)}, which makes the envelope's payload
bytes visible to any subsequent acquire on the same cluster.

\item \textbf{Consumer acquire}: at the start of \textsc{Deliver},
the consumer issues a matching
\texttt{cuda::\allowbreak atomic\_thread\_fence(\allowbreak
acquire,\ scope::cluster)} before reading the envelope payload.

\item \textbf{Cursor atomicity}: the producer and consumer cursors
of each queue are 32-bit values updated with
\texttt{cuda::atomic\_ref<uint32\_t>}-style operations.  Per the
CUDA C++ programming
guide~\citep{nvidia2023cooperative,nvidia2023hopperguide}, the
built-in \texttt{atomicAdd} and \texttt{atomicCAS} functions are
atomic only at their specified scope and behave as
\texttt{memory\_order\_relaxed} with respect to the surrounding
ordering.  The runtime therefore pairs each cursor update with an
explicit scoped fence (the release/acquire pair above) rather than
relying on the atomic operation alone for ordering; cursor atomicity
alone gives single-value atomicity, not system-scope sequential
consistency.
\end{enumerate}

The first two are the standard ``release-acquire'' synchronization
pattern; together they establish a happens-before edge from the
\textsc{SendOK} of an envelope to the matching
\textsc{Deliver}~\citep{lustig2014abstraction}.  The third gives the
queue's bounded-buffer property: the producer never overwrites a
slot the consumer has not yet released, and vice versa.

The formal model elides the fences: it treats \textsc{SendOK} and
\textsc{Deliver} as atomic single-step events.  This is a sound
abstraction provided the runtime emits the fences described above;
without them the K2K calculus would admit traces in which the
delivered envelope's payload bytes were partially observable.  The
runtime's K2K implementation, in
\texttt{crates/\allowbreak ringkernel-\allowbreak cuda/\allowbreak
src/k2k.cu}, emits the fences at the right places; we sketch the
correspondence in \Cref{sec:implementation:k2k}.

\subsection{Snapshot Atomicity}
\label{sec:memory:snapshot}

The snapshot--restart protocol of \Cref{sec:supervision:restart}
requires that \(\mathit{snap}\) is atomically equal to
\(\mathit{live}\) at the moment of \textsc{Activate}.  The runtime
realizes this by serializing \textsc{Activate} against
\textsc{Process}: the supervisor pauses the actor's execution
(transitioning \Active{} to \Initializing{} via a flag in mapped
memory polled at the start of every \textsc{Process} iteration),
copies \(\mathit{live}\) to \(\mathit{snap}\) using
\texttt{cudaMemcpyAsync} on a supervisor-private stream, awaits
the copy's completion via \texttt{cudaEventSynchronize}, then sets
the flag back to \Active{}.  The full sequence happens within a
single host-side critical section under a \texttt{Mutex} held by
the supervisor; concurrent supervisor actions on the same actor are
serialized by that mutex.

The snapshot copy is consistent with the device's memory consistency
model along a specific synchronization path.  \texttt{cudaMemcpyAsync}
itself is not a formal release-acquire primitive in the sense of
\Cref{sec:memory:consistency}; per the CUDA asynchronous-execution
documentation~\citep{nvidia2023cooperative} it is a
\emph{stream-ordered} operation whose completion is observable via
\texttt{cudaEventRecord}/\texttt{cudaEventSynchronize} or
\texttt{cudaStreamSynchronize}.  What establishes the happens-before
edge we rely on is the full sequence
(i)~device-side flag flip \Active{}\,\(\to\)\,\Initializing{}
committed via an explicit release fence,
(ii)~\texttt{cudaMemcpyAsync} queued on the supervisor's stream,
(iii)~\texttt{cudaEventRecord} on the same stream, and
(iv)~host-side \texttt{cudaEventSynchronize} before the flag is
flipped back.  The combination, not any single step, gives the
formal guarantee the model needs.  The formal model treats the
composite as instantaneous; the runtime's actual copy cost is
bounded by HBM bandwidth (\(\sim\)3~TB/s on H100, so a 4-MiB
snapshot copies in \(\sim\)1.4~$\mu$s).
